\documentclass{article}
\usepackage{amsmath}
\usepackage[hidelinks]{hyperref}
\title{Shared stage debits, qualified cache reads and a consuming cut}
\date{}
\begin{document}
\maketitle

\paragraph{Target and pins.}
Can prior source and transport events generate a nonrectangular family
of observable evidence, with a consequence for a guarded live recovery?
This is a held-in realization of the comparison docket, not
a claim that a production stack has this path topology or achieves the
latency calculated here. The replay executes pinned RWARE movement and
grid update code and the selected pinned SubT dispatch, radio and
application callback code; exact Git blobs are in \verb|SOURCE.tsv|.
Endpoint registration, causal forwarding, stage layout, external link
seam, caches and acquisition/consumption rules are authored. The
radio seam permits only an unambiguous receive or a failed path; no
noise, propagation or field link-quality assumptions are inferred.
The simulator's local motion order and transport clock are synchronized
by this fixture; there is no distributed simultaneity theorem.

\paragraph{Construction.}
One carrier's real movement produces old reads of occupancy of a named
bay. Each read is taken once from the current grid after the appropriate
move or stay. Either bit can recur, and the construction realizes all
six-bit old cache vectors without inventing a false sensor value.
Then the source latches a new occupancy bit at the protected cut. The
six labeled copies travel along logical forwarding paths
\[
(0,4),\quad(1,3,4),\quad(1,2),\quad(0,1),\quad(2,4),\quad(0,2,3)
\]
in label order. Only a real callback advances to the next hop or
updates the addressed cache. Each stage has a separate sender and
recipient, capacity four bytes and at most three one-byte refreshes.
One unreceived four-byte background send at a selected stage is real
TX debt but does not create a receive. At quarter-tick four its age
is exactly one radio second, and the native capacity test still
includes it. At tick five it is outside the native window. Other
stages have headroom. Thus the blocked paths are exactly those that
visit a live depleted stage, or the union if two such sends remain
live. The second cause is an actual packet debit, not a second mask
applied to caches. Actual unrelated radio delivery and carrier motion
occur in every fixture. The separate source bit is hidden from the
controller during acquisition and inspected by the evaluator only.

\paragraph{Deciding under a protected snapshot.}
For each live-debit mode let $F$ be its actual stale path set. A bare
cache word is legal for bit $b$ iff it agrees with $b$ outside $F$
for some admitted mode. This follows in both directions: fresh
callbacks copy the latch unchanged; each cache value left by a
failed refresh is realizable by a past native read. On a fixed set
of queried labels $V$, opposite truths have a common view iff
some pair of possible stale sets $F,G$ satisfies $V\subseteq F\cup G$.
If no such pair exists, some queried label would be simultaneously
fresh in the two candidate explanations. Conversely, assign the
common fresh-required values outside the stale sets and arbitrary
values where both may be stale. These choices are consistent exactly
when the queried labels are in the union.

At the inclusive cut with at most one depleted stage this rules out
every five-label fixed plan but all six suffice. For adaptive reads,
the replay computes the exact decision tree: a view with one possible
true bit costs zero; a mixed view branches on any unread label,
taking the maximum child cost plus one, then minimizing over labels.
Conditioning never uses the hidden debit index. The optimum is
five. With two depleted stages, even all six bare values can be
identical for opposite truth. For example, the old bits can leave the
word $000001$ in either case: only label \texttt{e} is fresh for
truth zero using stages $0,1$; only label \texttt{f} is fresh for
truth one using stages $1,4$. The replay computes the witnesses
from actual callback traces.

A qualified read returns the value and whether this cache received
the specific latched sample, not a generic link-up indication, message
age inferred from arrival, or guarantee of current bay state. In
this fixture the additional bit is local metadata derived from a
real final callback at the one labeled latch. A reused or forged
qualification is outside the model. With one depleted stage the
fixed pair $\{\texttt{a},\texttt{c}\}$ always includes a fresh cache;
every individual label has a possible failed refresh, so two
worst-case reads are necessary even with adaptivity. With two live
debits, a fresh label still exists in every mode, but the replay
finds worst-case depth six with either fixed or adaptive
acquisition when the read contains the freshness bit.
After expiry one read suffices in either representation.

\paragraph{Consumption and comparison.}
All 4,096 generation cases pass. At the inclusive cut, no carrier
reentry is allowed between the latch and the decision's actual
application callback. The fixture uses the current bit to select
direct when clear and a distinct logical bypass when occupied;
only these authorized callback effects are verified, not physical
bypass motion. One serialized cache query is charged one outer
quarter-tick. The supported due time is tick nine after the tick-four
latch, so the optimal adaptive bare and fixed qualified acquisition
complete all 768 single-debit-or-none executions; the exact fixed
bare plan cannot guarantee service within that budget. The qualified
arm takes at most two and the bare adaptive arm at most five.
A policy with stronger observability (e.g. one atomic direct reader,
or a trusted live stage ledger) is not dominated by the opaque-cache
comparison. The explicit-world solver itself attains the optimum;
there is no gap against it. A separate native reentry witness violates
the hold after an all-fresh acquisition: the same qualified decision
can still generate a real radio callback despite no longer being
lawful. The lost assumption is occupancy stability between the two
cuts, not sensor accuracy.

These mechanisms generate actual stale combinations for a specified
architecture; they do not derive this architecture's scope from the
inventory labels or certify what faults an external mission admits.
The six paths were selected to realize the earlier discriminator.
No hardware, physically modeled bypass, independent transfer,
confirmatory benchmark, runtime timing measurement or I1 acceptance
is claimed.
\end{document}
